\begin{lemma}[Extending a core by exceptional rows]
Let $H,F$ be finite simple graphs on the same vertex set, and let
$U,W$ be disjoint finite sets with $U_0\subseteq U,W_0\subseteq W$.
Suppose $F\subseteq H$ and every edge of $F$ joins $U_0$ to $W_0$.
Let $k$ be a nonnegative integer. Suppose each vertex of $U\setminus U_0$
has at least $k$ neighbors in $W_0$ in $H$, and each vertex of
$W\setminus W_0$ has at least $k$ neighbors in $U_0$ in $H$.
There is a simple graph $F\subseteq G\subseteq H$ on the same vertex set
which agrees with $F$ between $U_0,W_0$, gives each exceptional vertex
exactly $k$ neighbors in the opposite full part, and satisfies
\[
\deg_F(u,W_0)\leq\deg_G(u,W)\leq\deg_F(u,W_0)+|W\setminus W_0|
\quad(u\in U_0),
\]
and the analogous inequalities with $U,W$ exchanged.
\end{lemma}
\begin{proof}
For each exceptional vertex choose a $k$-element subset of its neighbors
in the opposite core. Such a subset exists by the cardinality hypothesis.
Add all chosen edges, in both orientations, to $F$ to obtain $G$.
All added edges belong to $H$, which proves containment and looplessness;
including both orientations proves symmetry. Finite adjacency is decidable.
An added edge has exactly one exceptional endpoint, since the other endpoint
is in the opposite core and the full parts are disjoint. In particular no
added edge joins two core vertices, so $G$ agrees with $F$ on the core.
The support hypothesis on $F$ gives an exceptional vertex no old neighbors.
No other exceptional vertex can prescribe an edge incident to it, because
every prescribed opposite endpoint is in a core. Its opposite-part neighbor
set is therefore exactly its chosen set, of size $k$.
For a core vertex in $U_0$, all old neighbors lie in $W_0$, all such edges
are retained, and every new opposite-part neighbor lies in $W\setminus W_0$.
Its neighbor set is thus contained in the union of its old neighbor set
with $W\setminus W_0$. Taking cardinalities proves the two bounds.
The identical argument for a vertex in $W_0$ proves the other bounds.
\end{proof}
